%
%  chapter2.tex — Background and State of the Art
%
\chapter{Background and State of the Art}
\label{cha:background}

\section{E-Learning Platforms}
\label{sec:elearning}

The e-learning industry has grown substantially over the past decade. Platforms such as Coursera, edX, Udemy, and Khan Academy have demonstrated that digital education can reach millions of learners globally. However, their core model — video lectures supplemented by multiple-choice assessments — has been repeatedly criticized for low engagement and poor knowledge retention~\cite{jordan2014initial}.

Duolingo~\cite{vesselinov2012duolingo} is a notable exception: by applying gamification rigorously to language learning, it achieves significantly higher daily active usage than traditional platforms. Play4Change draws inspiration from this approach while extending it to sustainability and digital literacy topics and adding AI-driven content generation.

\section{Gamification in Education}
\label{sec:gamification}

Gamification is defined by Deterding et al.~\cite{deterding2011game} as the use of game design elements in non-game contexts. In education, the most effective mechanics include:

\begin{itemize}
    \item \textbf{Daily streaks}: reinforce consistent engagement through loss aversion.
    \item \textbf{Leaderboards}: provide social proof and healthy competition.
    \item \textbf{Achievement badges}: signal mastery and milestone completion.
    \item \textbf{Peer review}: builds critical thinking and collaborative knowledge through distributed evaluation.
\end{itemize}

Play4Change incorporates all of these mechanics. Gamification must serve learning, not replace it — badges are awarded for genuine comprehension, not mere participation.

\section{Adaptive Learning and Flow State}
\label{sec:adaptive}

Adaptive learning systems adjust the difficulty, pacing, and content of lessons based on each learner's performance~\cite{vanlehn2011relative}. The goal is to maintain learners in the \emph{flow state}~\cite{csikszentmihalyi1990flow} — the optimal cognitive zone between boredom and anxiety.

Play4Change implements adaptive difficulty through a struggle detection mechanism (see Section~\ref{sec:struggle_detect}) that triggers AI-generated remediation branches when a learner fails the same task twice consecutively.

\section{Generative AI and Retrieval-Augmented Generation}
\label{sec:genai}

Large Language Models (LLMs) have transformed content generation. However, raw LLM output suffers from hallucination and lack of grounding in specific source material. Retrieval-Augmented Generation (RAG)~\cite{lewis2020retrieval} addresses this by augmenting generation with a retrieval step: documents are embedded into a vector database and semantically similar passages are retrieved at inference time to ground the model's response.

Play4Change uses the \textbf{Mistral} LLM~\cite{mistral2024} for content generation, combined with a \textbf{pgvector} index over educator-supplied materials. Mistral is selected for its strong performance-to-cost ratio and self-hostable deployment option. \textbf{LangChain4j}~\cite{langchain4j2024} provides the JVM integration layer between the Spring Boot server and Mistral, behind a provider-agnostic interface (see Section~\ref{sec:ai_module}).

\section{Peer Review as an Assessment Mechanism}
\label{sec:peer_review_bg}

Peer review is a pedagogically superior assessment mechanism for applied, open-ended tasks. The act of evaluating another learner's work is itself a form of retrieval practice — one of the most evidence-backed learning techniques~\cite{roediger2011retrieval}. Calibrated peer review in higher education has demonstrated that distributed assessment across three independent reviewers produces robust signals even for subjective tasks~\cite{reily2010two}.

Play4Change uses majority-vote peer review (2 of 3 verdicts) for \emph{todo-action} tasks where learners photograph real-world applications of their learning. This achieves zero marginal AI cost per submission while simultaneously creating learning events for the reviewers.

\section{Sustainability Education and the UN SDGs}
\label{sec:sdgs}

The United Nations' 2030 Agenda defines 17 Sustainable Development Goals~\cite{unsdgs2015} covering poverty, health, education, climate action, and more. Despite global commitments, integrating SDG content into daily learning routines remains a challenge. Play4Change curates topics around the SDGs and digital citizenship, enabling educators to contribute materials that the AI automatically structures into SDG-tagged learning paths.

\section{Passwordless Authentication}
\label{sec:passwordless}

Traditional password-based authentication is a leading cause of security breaches. The FIDO2 standard~\cite{fido2} and related approaches (magic links, OTPs, social OAuth) eliminate shared secrets entirely. Play4Change adopts passwordless authentication as a first principle: magic links delivered via email for users without social accounts, and OAuth 2.0 ID token verification (Google, Facebook) for users who prefer social login. No password hashes are stored anywhere in the system (see Section~\ref{sec:security}).

\section{Cross-Platform Mobile Development}
\label{sec:crossplatform}

Three major approaches compete for cross-platform mobile development:

\begin{itemize}
    \item \textbf{React Native} (Meta): JavaScript/TypeScript bridge to native components.
    \item \textbf{Flutter} (Google): Dart with a custom rendering engine.
    \item \textbf{Kotlin Multiplatform (KMP)}~\cite{kotlinmultiplatform2024}: shares business logic and data layers across platforms while keeping UI native.
\end{itemize}

Play4Change uses KMP because it allows the team to share domain logic, network code, serialization, and database queries across iOS and Android, while writing native UI in Jetpack Compose (Android) and SwiftUI (iOS). \textbf{Decompose}~\cite{decompose2024} is used for navigation and lifecycle management within KMP, providing a component-based navigation model that works identically on both platforms without relying on Android-specific ViewModels.

\section{Typed Error Handling: Railway Oriented Programming}
\label{sec:rop}

Railway Oriented Programming~\cite{wlaschin2014rop} is a functional design pattern for building pipelines where each step either continues on the happy path or short-circuits with a typed error. It avoids thrown exceptions for \emph{expected} domain failures — validations, not-found cases, business rule violations — replacing them with explicit return types.

Play4Change uses \textbf{Arrow}~\cite{arrowkt2024} \texttt{Either<E, A>} and the \texttt{Raise<E>} DSL for this purpose. The custom error hierarchy (\texttt{AppError} and its sealed subclasses) is domain-specific and Arrow-independent; Arrow provides the combinators (\texttt{bind()}, \texttt{ensure()}, \texttt{mapError()}, \texttt{zip()}) (see Appendix~\ref{app:adrs}, ADR-001).

\section{Backend Technologies}
\label{sec:backend}

Spring Boot~\cite{springboot2024} with Kotlin is the chosen backend framework. Kotlin's null-safety, coroutine support, and seamless interoperability with the Java ecosystem make it an ideal fit. The backend is containerized with Docker for reproducible deployments. \textbf{MinIO}~\cite{minio2024} provides S3-compatible object storage for educator-uploaded PDF materials, running as a Docker Compose service alongside the application.

\section{Architectural Patterns}
\label{sec:architecture_patterns}

\subsection{Domain-Driven Design}
\label{ssec:ddd}

Domain-Driven Design (DDD)~\cite{evans2003domain} organizes software around the business domain rather than technical concerns. Key DDD building blocks used in Play4Change include \emph{aggregates}, \emph{value objects}, \emph{domain events}, and \emph{bounded contexts} — each with its own model and repository.

\subsection{Hexagonal Architecture}
\label{ssec:hexagonal}

The hexagonal architecture~\cite{cockburn2005hexagonal} separates the application core from its infrastructure. The core defines \emph{ports} (interfaces), and \emph{adapters} implement those ports for specific technologies. This enables thorough unit testing of domain logic without any infrastructure.

\subsection{Clean Architecture}
\label{ssec:clean}

Clean Architecture~\cite{martin2017clean} enforces a strict dependency rule: inner layers know nothing about outer layers. In Play4Change: \textbf{Domain} → \textbf{Application} → \textbf{Infrastructure} → \textbf{Presentation}. Domain logic is framework-agnostic and testable in isolation.

\section{Observability: Metrics and Monitoring}
\label{sec:observability_bg}

Production systems require visibility into runtime behaviour. \textbf{Micrometer}~\cite{micrometer2024} provides a metrics façade that decouples instrumentation from any specific backend. \textbf{Prometheus}~\cite{prometheus2024} scrapes Micrometer's HTTP endpoint and stores time-series data. \textbf{Grafana}~\cite{grafana2024} queries Prometheus with PromQL and renders dashboards. This stack runs entirely within Docker Compose, making the full observability pipeline reproducible from source.

\section{Related Work and Gap Analysis}
\label{sec:related}

Table~\ref{tab:related_work} summarizes the key differentiators of Play4Change relative to existing platforms.

\begin{table}[ht]
    \caption{Comparison of Play4Change with related platforms.}
    \label{tab:related_work}
    \vskip0.5em
    \centering
    \begin{tabular}{lcccc}
    \toprule
    Feature & Duolingo & Coursera & Khan Academy & Play4Change \\
    \midrule
    AI content generation    & \texttimes & \texttimes & \texttimes & \checkmark \\
    Adaptive difficulty      & \checkmark & \texttimes & Partial    & \checkmark \\
    Daily streak mechanics   & \checkmark & \texttimes & \texttimes & \checkmark \\
    Peer review tasks        & \texttimes & Partial    & \texttimes & \checkmark \\
    SDG-focused topics       & \texttimes & Partial    & Partial    & \checkmark \\
    Passwordless auth        & \texttimes & \texttimes & \texttimes & \checkmark \\
    KMP cross-platform       & \texttimes & \texttimes & \texttimes & \checkmark \\
    Open-source stack        & \texttimes & \texttimes & \texttimes & \checkmark \\
    \bottomrule
    \end{tabular}
\end{table}
